\chapter{“보냈어”와 “써도 돼”}

모두의 공책에는 여러 종류의 줄이 적혀요.
이 책에서는 그중 딱 두 가지만 기억하면 돼요.

\section{첫 번째 줄: “보냈어”}

\begin{center}
\fbox{\parbox{0.8\linewidth}{\centering 3월 4일 \quad 지우가 하준이에게 별 스티커 2장을 보냈다.}}
\end{center}

이런 줄을 \textbf{송금}이라고 해요.
스티커가 한 통에서 다른 통으로 실제로 옮겨 간 기록이에요.
송금은 아주 자주 일어나요.
물건값을 낼 때, 선물을 줄 때, 내 통에서 내 다른 통으로 옮길 때,
그리고 로봇들이 스티커를 이리저리 주고받을 때도 송금 줄이 생겨요.

\section{두 번째 줄: “써도 돼”}

\begin{center}
\fbox{\parbox{0.8\linewidth}{\centering 3월 4일 \quad 지우가 문구점 로봇에게\\“내 통에서 파란 스티커를 10장까지 꺼내 가도 돼”라고 허락했다.}}
\end{center}

이런 줄을 \textbf{허락}이라고 해요.
이번에는 스티커가 하나도 움직이지 않았어요.
지우는 문구점 로봇에게 자기 통의 문을 조금 열어 준 거예요.

왜 이런 허락이 필요할까요?
문구점 로봇은 규칙대로만 움직이는 프로그램이에요.
지우가 연필을 사려고 할 때마다 스티커를 직접 보내는 대신,
“10장까지는 네가 알아서 꺼내 가”라고 미리 허락해 두면 로봇이 필요한 만큼만 꺼내 가요.
블록체인의 앱들은 대부분 이렇게 움직여요.
교환 로봇에게 스티커를 바꿔 달라고 하려면, 먼저 교환 로봇에게 허락을 해 줘야 하지요.

허락한 양을 \textbf{허락 한도}라고 불러요.
허락은 나중에 바꿀 수 있어요.
“이제 5장까지만”이라고 새로 적을 수도 있고,
“이제 0장”이라고 적어서 허락을 거둘 수도 있어요.
공책에 허락 줄이 여러 번 적혀 있으면 \emph{가장 마지막 줄}이 지금의 허락이에요.

\begin{figure}[h]
\centering
\begin{tikzpicture}[box/.style={draw=inkblue, fill=skyfill, rounded corners=1.5mm, minimum width=17mm, minimum height=9mm, font=\small}]
  \node[font=\small\bfseries, color=sendcol, anchor=west] at (-0.9,0.95) {보냈어 (송금)};
  \node[box] (j1) at (0,0) {지우 통};
  \node[box] (h1) at (5.6,0) {하준 통};
  \draw[send] (j1) -- node[above, font=\scriptsize, color=black] {별 스티커 2장이} node[below, font=\scriptsize, color=black] {실제로 옮겨 감} (h1);
  \node[font=\small\bfseries, color=allowcol, anchor=west] at (-0.9,-1.45) {써도 돼 (허락)};
  \node[box] (j2) at (0,-2.4) {지우 통};
  \node[robot] (r2) at (5.6,-2.4) {문구점 로봇};
  \draw[allow, dashed] (j2) -- node[above, font=\scriptsize, color=black] {파란 스티커} node[below, font=\scriptsize, color=black] {10장까지 꺼내 가도 돼} (r2);
\end{tikzpicture}
\caption{송금은 스티커가 실제로 옮겨 가는 일이고, 허락은 통의 문을 열어 두는 일이에요.}
\end{figure}

\section{허락에는 믿음이 담겨 있어요}

자, 이제 이 책에서 가장 중요한 생각이 나와요.
천천히 읽어 주세요.

송금은 꼭 상대를 믿어서 하는 일이 아니에요.
처음 본 가게에 물건값을 낼 때도 송금을 하고,
그냥 스티커를 이 통에서 저 통으로 옮길 때도 송금을 해요.
송금 한 번은 “한 번 주고받았다”는 뜻일 뿐이에요.

그런데 허락은 달라요.
허락은 내 통의 문을 열어 두고 \emph{계속} 열어 두는 일이에요.
상대가 언제든 들어와서 정해진 만큼 꺼내 갈 수 있어요.
여러분은 오늘 처음 만난 사람에게 집 열쇠를 맡기지 않을 거예요.
문을 열어 둔다는 건 “나는 너를 믿어”라고 말하는 것과 같아요.

어려운 말로는 이런 믿음의 표시를 \textbf{보증}이라고 해요.
제 논문의 출발점은 바로 이 생각이에요.

\begin{center}
\begin{tcolorbox}[colback=white, colframe=allowcol, boxrule=0.8pt, arc=2mm, width=0.9\linewidth, halign=center]
\titlefont\bfseries 허락 줄은 “나는 너를 믿어”라는 기록이다.\\그렇다면 허락 줄로 평판 점수를 만들 수 있지 않을까?
\end{tcolorbox}
\end{center}

\section{허락을 받는 쪽은 대부분 로봇이에요}

한 가지 알아 둘 것이 있어요.
허락을 \emph{받는} 쪽은 대부분 사람이 아니라 로봇이에요.
문구점 로봇, 교환 로봇, 대출 로봇처럼요.
나중에 진짜 기록을 살펴보니, 허락을 받아 둔 통 1,335개 가운데 1,174개가 로봇이었어요.
그러니 허락 줄로 만든 점수는 주로 “사람들이 어떤 앱을 믿는가”를 말해 줘요.
그리고 사람의 통은 주로 허락을 \emph{해 주는} 쪽에서 나타나요.

\section{허락 줄은 송금 줄보다 훨씬 적어요}

또 하나 눈여겨볼 점이 있어요.
허락은 송금보다 훨씬 드물어요.
제가 살펴본 지갑들의 여섯 달 치 기록에서 허락 줄은 약 45만 줄이었는데, 송금 줄은 약 280만 줄이었어요.
여섯 배쯤 차이가 나지요.
이 차이가 12장에서 계산 속도의 차이로 이어져요.

\begin{deeper}[진짜 용어]
\begin{itemize}
\item 공책에서 송금 줄은 “송금 사건”, 허락 줄은 “승인 사건”이라는 이름으로 적혀요. 진짜 기록에서는 영어 이름 \texttt{Transfer}와 \texttt{Approval}로 찾을 수 있어요.
\item 허락 한도는 “허용량”이라고도 해요. 허락을 해 주는 쪽은 “주인”, 허락을 받는 쪽은 “사용자”라고 불러요.
\item 논문에서는 (토큰, 주인, 사용자)마다 가장 마지막 승인 사건의 값을 “최신 허락”으로 삼았어요. 값이 0이면 허락을 거둔 것으로 봐요.
\item 2025년 12월부터 2026년 5월까지 연구 대상 지갑과 관련된 줄은 승인 사건 453,545줄, 송금 사건 2,801,988줄이었어요.
\end{itemize}
\end{deeper}

\begin{think}
송금은 많이 하는데 허락은 거의 하지 않는 사람은 어떤 사람일까요?
반대로 허락을 아주 많은 로봇에게 해 둔 사람은 어떤 사람일까요?
\end{think}

\begin{learned}
\begin{itemize}
\item 송금은 스티커가 실제로 옮겨 간 기록이고, 허락은 통의 문을 열어 둔 기록이에요.
\item 허락에는 “나는 너를 믿어”라는 뜻이 담겨 있어요. 이것이 이 논문의 출발점이에요.
\item 허락은 대부분 로봇(앱)이 받고, 송금보다 훨씬 드물어요.
\end{itemize}
\end{learned}
